\section{Conclusion}
DiagPriv tests whether a release can hide selected literals and retain the
relations needed for RCA. The three synthetic families require linkage,
shape, or time. The paired temporal result fell from 3/4 discriminated pairs
with raw data to 2/4 with relative time and 0/4 with ablation. This small
result is consistent with a useful temporal dependency, not a general causal
estimate.

Literal removal can still expose linkage, shape, cardinality, time, and
context. The paper therefore reports RCA, evidence support, and exposure
separately. Future work needs more incidents, models, human baselines, and a
frozen held-out test before it can evaluate a task-aware selection policy.
